Contemporary neural text generators enable rapid composition of long-form documents, but they do not by themselves guarantee that produced assertions are verifiable or that supporting sources are cited.
Automated long-form generation systems can produce unsupported or inaccurate statements, invent citations, or omit references that a human author would provide.
For workflows that target scholarly or otherwise high-assurance outputs, these failure modes are unacceptable: authors need explicitly verifiable claims, machine-parseable evidence links, and a reproducible audit trail that supports downstream checking and revision.

This work is motivated by that practical need.
We investigate whether structuring the generation process around an explicit document graph improves the reliability, citation quality, and verifiability of section-level outputs.
Intuitively, a document graph makes section boundaries, dependencies, and evidence flow explicit, allowing retrieval and generation modules to operate with richer context than linear prompting alone.

Concretely, we represent a draft document as a directed graph G = (V,E) in which each node v $\in V$ corresponds to a section (or subsection) and edges e $\in E$ encode inter-section relations such as dependency, evidence-flow, refinement, or contrast.
Each node is annotated with a small schema of attributes: target length, rhetorical role (e.g., background, method, result), local goals or claims, and a placeholder for attached evidence pointers.
Edges carry typed metadata that indicate, for example, that section A supplies background evidence for section B, or that section C refines claims made in section D.

The document graph plays three operational roles in the proposed workflow.
First, it constrains generation ordering via a scheduler that respects dependencies: nodes are generated in a topologically consistent order, leaves may be produced first, and independent subgraphs can be parallelized when safe.
The scheduler also supports prioritized re-generation when upstream edits change evidence or claims.
Second, the graph informs retrieval by producing context-aware queries for each node: a node's query is constructed from its local goals plus summaries of upstream supporting nodes, enabling focused retrieval that takes inter-section context into account.
Third, the graph governs assembly by mapping structured section outputs into a coherent LaTeX manuscript: citation pointers are resolved into a consolidated refs.bib, citation identifier conflicts and cross-reference consistency are resolved deterministically, and provenance metadata is embedded to yield an auditable record.

The workflow itself proceeds in four stages.
Stage 1: an initial document graph is constructed from a high-level outline or template; human authors or automated outline parsers provide node headings and annotations such as intended claims and length targets.
Stage 2: for each node the system retrieves candidate evidence items from curated internal notes and a collection of small PDFs; retrieved items are ranked and attached to nodes as evidence pointers, along with retrieval scores and provenance identifiers.
Stage 3: leaf and scheduled sections are generated under constrained prompts that require JSON-compliant outputs: explicit fields include title, an array of abstracted claims, a sequence of paragraphs with span indices, citation slots that reference evidence pointer IDs and include span-level start/end indices, and provenance metadata that records retrieval identifiers and timestamps.
In this stage each factual claim must either be linked to at least one supporting evidence pointer or be explicitly labeled as speculative; generators may also return ranked evidence suggestions and confidence scores.
Stage 4: validated JSON section outputs are assembled into a single LaTeX manuscript and a unified refs.bib; the assembler deterministically canonicalizes citation keys, resolves duplicate references, replaces pointer IDs with bibkeys in \cite commands, and emits a machine-readable audit manifest that maps each asserted claim span to the supporting retrieval items.

For evaluation we ask: how does graph-guided sectioning affect coherence, citation coverage, and factual consistency relative to a baseline linear prompting pipeline?
We compare the graph-guided pipeline to a linear baseline across two topical domains.
Measured outcomes include citation coverage (fraction of non-speculative claims with attached evidence), hallucination rate (estimated via reference-based checks that verify whether cited sources actually support claimed facts), automated factual-consistency metrics where applicable, and blind reviewer assessments of coherence, usefulness, and verifiability.
To encourage conservative outputs, we enforce a generation regime in which section generators must attach supporting evidence for factual claims or mark them as speculative; low-confidence or low-recall retrievals are surfaced to authors rather than allowing unsupported assertions to propagate.

This paper makes three primary contributions.
First, we introduce a document-graph formulation for long-form generation that encodes section hierarchy, dependency types, and explicit evidence links.
Second, we present a section scheduler and evidence-injection mechanism that together produce JSON-constrained section outputs with fine-grained provenance.
Third, we describe an assembly procedure that converts section-level JSON into a LaTeX manuscript with a consolidated refs.bib and an embedded audit trail suitable for downstream verification and iterative revision.

We also acknowledge important sources of uncertainty and limits to generality.
Section headings and rhetorical roles do not uniformly indicate citation intent across disciplines, so section-aware citation signals can be noisy and require cautious interpretation \cite{web_thelwall_2019_should_citations_be_counted_separately_from_each}.
Retrieval modules that operate on flat corpora can miss fragmented, dispersed, or multi-hop evidence; prior work suggests that explicitly graph-structured retrieval can help identify and combine such evidence items \cite{web_mongiov_2021_graph_based_retrieval_for_claim_verification_ove}.
Finally, structured and cryptographically-aware evidence records are increasingly discussed for regulated audit settings; these approaches underscore the importance of durable, machine-verifiable provenance for high-assurance deployments and inform our audit-trail design choices \cite{web_kao_2025_quantum_adversary_resilient_evidence_structures_}.

In sum, this introduction motivates a graph-guided prompting workflow that prioritizes verifiability and auditability in long-form generation.
The remainder of the paper details the document-graph representation and scheduler (Section 3), the retrieval and constrained-generation procedures (Sections 3.4–3.5), experimental design (Section 4), empirical results (Section 5), and a conservative discussion of strengths, limitations, and future directions (Sections 6–7).